\paragrafos{Δεκαδικά κλάσματα και δεκαδικοί: ανάμεσα στους ακέραιους}{Α.3.1, «Δεκαδικά κλάσματα – Δεκαδικοί αριθμοί – Διάταξη – Στρογγυλοποίηση»}

\begin{anagke}
Αθλητική μέρα στο καινούργιο σχολείο, και για πρώτη φορά οι δύο γειτονιές τρέχουν στην ίδια πίστα. Η κυρία Αλεξάνδρα κρατάει το χρονόμετρο στα 60 μέτρα και η Πολυξένη γράφει τους χρόνους στο τετράδιο, γιατί στη Θεσσαλονίκη είχε τρέξει σε σύλλογο.

Ο Νίκος τερματίζει πρώτος και η καθηγήτρια φωνάζει: «Εννιά και ογδόντα!». Λίγο μετά τερματίζει ο Βασίλης: «Εννιά και εβδομήντα πέντε». Ο Βασίλης σκύβει το κεφάλι: «Το εβδομήντα πέντε είναι πολύ μεγαλύτερο από το οκτώ. Άρα έκανα πολύ περισσότερο χρόνο από σένα, έχασα με μεγάλη διαφορά». Ο Νίκος διαμαρτύρεται: «Μα έφτασες σχεδόν μαζί μου, με είχες πάνω από την πλάτη μου!». Η Πολυξένη γράφει τους δύο χρόνους, $9{,}8$ και $9{,}75$, και παρατηρεί ότι δεν είναι ίδιου μήκους: «Το ένα έχει ένα ψηφίο μετά την υποδιαστολή και το άλλο δύο. Πώς τα συγκρίνεις;»

Ο Στέλιος προσθέτει το δικό του πρόβλημα: «Με τι μονάδα μετράμε τέτοιους χρόνους; Το δευτερόλεπτο είναι πολύ χοντρή μονάδα. Αν ο χρόνος μου είναι ανάμεσα στο 9 και στο 10, δεν μπορώ να πω ούτε 9 ούτε 10». Η Μυρτώ το είχε πάθει και στο θρανίο: το πλάτος του ήταν ανάμεσα στα 62 και στα 63 εκατοστά.

Οι φυσικοί δεν φτάνουν. Χρειαζόμαστε αριθμούς \emph{ανάμεσα} στους φυσικούς, που να τους κόβουν σε δέκατα, εκατοστά, χιλιοστά, όσο λεπτά χρειάζεται η μέτρηση. Τέτοιους αριθμούς τους ξέρουμε ήδη από τα κλάσματα, και τώρα θα τους γράψουμε με τρόπο που συγκρίνεται εύκολα.
\end{anagke}

\begin{lexi}{δεκαδικός, υποδιαστολή}{δεκαδικός, υποδιαστολή}
Η λέξη \lexiword{δεκαδικός} έρχεται από το «δέκα»: ο αριθμός είναι γραμμένος με βάση το 10, και μετά την υποδιαστολή συνεχίζονται οι θέσεις με δέκατα, εκατοστά, χιλιοστά. Η \lexiword{υποδιαστολή} είναι το κόμμα που χωρίζει τον ακέραιο από το δεκαδικό μέρος· «υπό-διαστέλλω» σημαίνει «χωρίζω από κάτω», και παλιότερα ήταν στο ύψος της βάσης των ψηφίων. Το \lexiword{δεκαδικό κλάσμα} είναι το κλάσμα που έχει παρονομαστή δύναμη του 10: $10,\,100,\,1\,000,\,10\,000,\dots$.
\end{lexi}

\begin{metafrasi}
Τρία δέκατα. & $\dfrac{3}{10}=0{,}3$\\
Εκατόν τριάντα πέντε εκατοστά. & $\dfrac{135}{100}=1{,}35$\\
Δύο ακέραια και επτά χιλιοστά. & $2{,}007=2+\dfrac{7}{1\,000}$\\
Ένα κλάσμα $\dfrac{\kappa}{\nu}$ γίνεται δεκαδικός. & $\kappa:\nu$\\
Οι ίδιοι χρόνοι, με τα ίδια ψηφία. & $9{,}8=9{,}80$\\
Ποιος δεκαδικός είναι μεγαλύτερος; & συγκρίνω ακέραιο, δέκατα, εκατοστά, …\\
Προσεγγίζω στα εκατοστά. & $9{,}7468\approx 9{,}75$\\
\end{metafrasi}

\tmima{Η θεωρία}

\tmimasmall{Δεκαδικά κλάσματα και δεκαδικοί αριθμοί}

Τα κλάσματα $\frac{3}{10}$, $\frac{825}{100}$, $\frac{53}{1\,000}$ και $\frac{1004}{10\,000}$ έχουν παρονομαστές τους $10$, $10^2$, $10^3$ και $10^4$. Είναι δεκαδικά κλάσματα.

\begin{orismos}{δεκαδικό κλάσμα και δεκαδικός αριθμός}
Δεκαδικό κλάσμα λέγεται κάθε κλάσμα με παρονομαστή δύναμη του 10. Κάθε δεκαδικό κλάσμα γράφεται ως \textbf{δεκαδικός αριθμός} με τόσα δεκαδικά ψηφία όσα μηδενικά έχει ο παρονομαστής του:
$$\dfrac{3}{10}=0{,}3,\qquad \dfrac{825}{100}=8{,}25,\qquad \dfrac{53}{1\,000}=0{,}053,\qquad \dfrac{1004}{10\,000}=0{,}1004.$$
\end{orismos}

Για να πάω από το δεκαδικό στο κλάσμα, κάνω το αντίστροφο: γράφω στον αριθμητή όλα τα ψηφία χωρίς την υποδιαστολή και στον παρονομαστή το $1$ με τόσα μηδενικά όσα τα δεκαδικά ψηφία. Άρα $3{,}25=\frac{325}{100}$.

Για να μετατρέψω ένα \emph{οποιοδήποτε} κλάσμα $\frac{\kappa}{\nu}$ σε δεκαδικό, κάνω τη διαίρεση $\kappa:\nu$, βάζοντας υποδιαστολή και μηδενικά στον διαιρετέο όταν τελειώσει το ακέραιο μέρος. Αν η διαίρεση τελειώνει, ο δεκαδικός είναι ακριβής, π.χ. $\frac{3}{8}=0{,}375$. Αν δεν τελειώνει ποτέ, όπως στο $\frac{1}{3}=0{,}333\ldots$, τον γράφουμε κατά προσέγγιση (θα το ξαναδούμε στο κεφάλαιο των ρητών).

\tmimasmall{Οι τάξεις των δεκαδικών ψηφίων}

Δεξιά από την υποδιαστολή οι θέσεις μικραίνουν κατά δέκα, όπως αριστερά μεγαλώνουν:

\begin{center}\small
\renewcommand{\arraystretch}{1.4}\setlength{\tabcolsep}{3pt}
\begin{tabular}{@{}cccccccc@{}}
\toprule
εκατοντάδες & δεκάδες & μονάδες & \textbf{,} & δέκατα & εκατοστά & χιλιοστά & δεκάκις χιλιοστά\\
$100$ & $10$ & $1$ & & $\frac{1}{10}$ & $\frac{1}{100}$ & $\frac{1}{1\,000}$ & $\frac{1}{10\,000}$\\
\bottomrule
\end{tabular}
\end{center}

Το επόμενο ψηφίο λέγεται εκατοντάκις χιλιοστό ($\frac{1}{100\,000}$), μετά εκατομμυριοστό, και ούτω καθεξής. Ο $305{,}0408$ έχει δεκαδικό μέρος $0{,}0408$ και διαβάζεται «τριακόσια πέντε ακέραια και τετρακόσια οκτώ δεκάκις χιλιοστά», γιατί το δεκαδικό μέρος $0{,}0408=\frac{408}{10\,000}$. Μπορούμε να τον αναλύσουμε:
$$305{,}0408=3\cdot100+0\cdot10+5\cdot1+\dfrac{0}{10}+\dfrac{4}{100}+\dfrac{0}{1\,000}+\dfrac{8}{10\,000}.$$

Αν προσθέσω ή σβήσω μηδενικά στο τέλος του δεκαδικού μέρους, ο αριθμός δεν αλλάζει: $9{,}8=9{,}80=9{,}800$, γιατί $\frac{8}{10}=\frac{80}{100}=\frac{800}{1\,000}$. Δεν αλλάζει όμως αν βάλω μηδενικό \emph{πριν} από τα ψηφία του δεκαδικού μέρους: $0{,}5\neq 0{,}05$.

\tmimasmall{Μια ιστορική σημείωση: ο Στέβιν και ο Νάπιερ}

Οι Βαβυλώνιοι και οι αρχαίοι Έλληνες αστρονόμοι χρησιμοποιούσαν κλάσματα με παρονομαστή 60, και οι έμποροι έβρισκαν δύσκολα τους υπολογισμούς με κλάσματα. Ο Φλαμανδός μαθηματικός \textbf{Σιμόν Στέβιν} (1548--1620) κατάλαβε ότι οι υπολογισμοί γίνονται πολύ πιο εύκολα με κλάσματα που έχουν παρονομαστή δύναμη του 10. Το 1585 έγραψε ένα μικρό βιβλίο, τη «Δεκάτη» (\emph{La disme}), για τους τοκογλύφους, τους τοπογράφους και τους εμπόρους, όπου εισήγαγε τον δεκαδικό υπολογισμό. Ο Στέβιν έγραφε τον $34{,}25$ ως $34_{(0)}\,2_{(1)}\,5_{(2)}$: ο δείκτης έδειχνε πόσες θέσεις μετά τη μονάδα βρίσκεται το κάθε ψηφίο. Το κόμμα (ή τελεία) ανάμεσα στον ακέραιο και στο δεκαδικό μέρος το πρότεινε αργότερα ο Σκοτσέζος \textbf{Τζον Νάπιερ} (1550--1617), και η απλότητά του επικράτησε.

\tmimasmall{Σύγκριση δεκαδικών}

\begin{kanonas}{σύγκριση δεκαδικών}
\begin{enumerate}[leftmargin=1.6em,itemsep=0.3ex,topsep=2pt]
\item Συγκρίνω πρώτα τα \textbf{ακέραια μέρη}. Μεγαλύτερο ακέραιο μέρος σημαίνει μεγαλύτερο αριθμό.
\item Αν τα ακέραια μέρη είναι ίσα, συγκρίνω τα δέκατα, μετά τα εκατοστά, μετά τα χιλιοστά, και σταματάω στο πρώτο ψηφίο που διαφέρει.
\item Για να μην μπερδεύομαι, γράφω στον αριθμό με τα λιγότερα δεκαδικά ψηφία μηδενικά στο τέλος, ώστε να έχουν το ίδιο πλήθος.
\end{enumerate}
\end{kanonas}

Έτσι ο χρόνος του Νίκου $9{,}8=9{,}80$ και του Βασίλη $9{,}75$: τα ακέραια ίσα, τα δέκατα $8>7$, άρα $9{,}80>9{,}75$. Ο Βασίλης συνέκρινε τα δεκαδικά μέρη $75$ και $8$ σαν να ήταν φυσικοί αριθμοί. Στην πραγματικότητα έκανε \emph{μικρότερο} χρόνο και κέρδισε, με διαφορά πέντε εκατοστά του δευτερολέπτου.

\tmimasmall{Στρογγυλοποίηση δεκαδικών}

Όταν δεν χρειαζόμαστε όλα τα ψηφία, τα κόβουμε και κρατάμε μια προσέγγιση. Για να στρογγυλοποιήσω σε μια δεδομένη τάξη (π.χ. στα εκατοστά):
\begin{itemize}[leftmargin=1.2em,itemsep=0.3ex,topsep=2pt]
\item κοιτάζω το \textbf{επόμενο} ψηφίο (στα χιλιοστά, αν στρογγυλοποιώ στα εκατοστά)·
\item αν είναι $0,1,2,3,4$ αφήνω το ψηφίο της τάξης όπως είναι και κόβω τα υπόλοιπα·
\item αν είναι $5,6,7,8,9$ αυξάνω το ψηφίο της τάξης κατά $1$ και κόβω τα υπόλοιπα.
\end{itemize}

\tmimasmall{Οι δεκαδικοί στην ευθεία}

Για να βάλω δεκαδικούς στην ευθεία των αριθμών, χωρίζω κάθε μονάδα σε $10$ ίσα μέρη (δέκατα) και, αν χρειάζεται, κάθε δέκατο σε $10$ ίσα μέρη (εκατοστά), και πάει λέγοντας. Δεν υπάρχει όριο: ανάμεσα σε δύο οποιουσδήποτε διαφορετικούς δεκαδικούς υπάρχει πάντα άλλος δεκαδικός, π.χ. ο μέσος τους, και άρα υπάρχουν άπειροι.

\begin{center}
\begin{tikzpicture}[x=1cm,y=1cm]
\draw[thick,-{Stealth}] (-0.3,0)--(10.7,0);
\foreach \i in {0,...,10}{\draw (\i,-0.12)--(\i,0.12);}
\node[below] at (0,-0.15) {$2$};
\node[below] at (5,-0.15) {$2{,}5$};
\node[below] at (10,-0.15) {$3$};
\foreach \i/\t in {1/2{,}1,3/2{,}3,8/2{,}8}{\node[below,font=\small] at (\i,-0.15) {$\t$};}
\fill[gblue] (7.5,0) circle (2.5pt);
\node[above,gblue] at (7.5,0.2) {$2{,}75$};
\fill[gblue] (3,0) circle (2.5pt);
\node[above,gblue] at (3,0.2) {$2{,}3$};
\end{tikzpicture}
\end{center}

Το $2{,}75$ βρίσκεται στη μέση του διαστήματος από το $2{,}7$ ως το $2{,}8$, που αντιστοιχεί στο σημείο $7{,}5$ των δέκατων.

\begin{prosoxi}
\begin{itemize}[leftmargin=1.2em,itemsep=0.3ex,topsep=2pt]
\item Δεν συγκρίνουμε τα δεκαδικά μέρη \emph{σαν φυσικούς}. Το $0{,}9$ είναι μεγαλύτερο από το $0{,}12$, γιατί $0{,}90>0{,}12$, αν και το $9<12$.
\item Το $0{,}50$ ισούται με το $0{,}5$, γιατί το μηδενικό στο τέλος δεν αλλάζει τίποτα, αλλά το $0{,}05$ είναι δέκα φορές μικρότερο.
\item Στρογγυλοποιούμε \emph{μια φορά}, από τον αρχικό αριθμό. Αν στρογγυλοποιήσουμε πρώτα στα εκατοστά και μετά στα δέκατα, μπορεί να βγει λάθος: $2{,}449\approx2{,}4$, αλλά $2{,}449\to2{,}45\to2{,}5$.
\item Το $\frac{2}{5}$ δεν έχει παρονομαστή δύναμη του 10, δεν είναι δηλαδή δεκαδικό κλάσμα, αλλά είναι ισοδύναμο με το δεκαδικό κλάσμα $\frac{4}{10}=0{,}4$.
\end{itemize}
\end{prosoxi}

\tmima{Λυμένα παραδείγματα}

\paradeigma{Ποιος κέρδισε;}
\tixero Οι χρόνοι στα 60 μ. του Νίκου και του Βασίλη είναι $9{,}8$ s και $9{,}75$ s.\par
\tithelo Ποιος έκανε τον μικρότερο χρόνο.\par
\possindeo Ισοσκελίζω τα δεκαδικά ψηφία: $9{,}8=9{,}80$. Τα ακέραια είναι ίσα (9) και στα δέκατα $8>7$, άρα $9{,}80>9{,}75$.\par
\apantisi Ο Βασίλης έκανε μικρότερο χρόνο ($9{,}75$ s) και κέρδισε με διαφορά πέντε εκατοστά του δευτερολέπτου.\par
\elegxos Σε εκατοστά, $980$ εκατοστά και $975$ εκατοστά: το $975$ είναι μικρότερο. Συμφωνεί.

\paradeigma{Από κλάσμα σε δεκαδικό και ανάποδα}
\tixero Τα κλάσματα $\frac{3}{8}$ και $\frac{7}{20}$ και ο δεκαδικός $0{,}625$.\par
\tithelo Να τα γράψω στην άλλη μορφή.\par
\possindeo Για το $\frac{3}{8}$ κάνω τη διαίρεση $3:8$: $3=3{,}000$ και $3000:8=375$, άρα $0{,}375$. Το $\frac{7}{20}$ γίνεται δεκαδικό κλάσμα αν πολλαπλασιάσω αριθμητή και παρονομαστή με $5$: $\frac{35}{100}=0{,}35$. Και ο $0{,}625$ έχει τρία δεκαδικά ψηφία, άρα $\frac{625}{1\,000}$, που απλοποιείται με το $125$ σε $\frac{5}{8}$.\par
\apantisi $\frac{3}{8}=0{,}375$, $\frac{7}{20}=0{,}35$ και $0{,}625=\frac{5}{8}$.\par
\elegxos $0{,}375\cdot 8=3$ και $\frac{5}{8}:1=0{,}625$ (αφού $5:8=0{,}625$). Σωστά.

\paradeigma{Διαβάζω και αναλύω}
\tixero Ο αριθμός $305{,}0408$.\par
\tithelo Να πω τι ψηφίο βρίσκεται σε κάθε θέση και να τον διαβάσω.\par
\possindeo Πριν από την υποδιαστολή: $3$ εκατοντάδες, $0$ δεκάδες, $5$ μονάδες. Μετά: $0$ δέκατα, $4$ εκατοστά, $0$ χιλιοστά, $8$ δεκάκις χιλιοστά. Το δεκαδικό μέρος είναι $0{,}0408=\frac{408}{10\,000}$.\par
\apantisi «Τριακόσια πέντε ακέραια και τετρακόσια οκτώ δεκάκις χιλιοστά».\par
\elegxos $\frac{4}{100}+\frac{8}{10\,000}=\frac{400}{10\,000}+\frac{8}{10\,000}=\frac{408}{10\,000}$. Σωστά.

\paradeigma{Στρογγυλοποιώ σε τάξεις}
\tixero Ο αριθμός $16{,}4963$.\par
\tithelo Να τον στρογγυλοποιήσω στην μονάδα, στα δέκατα, στα εκατοστά και στα χιλιοστά.\par
\possindeo Κοιτάζω κάθε φορά το επόμενο ψηφίο. Μονάδα: επόμενο το $4$ (<5) άρα $16$. Δέκατα: επόμενο το $9$ (≥5) άρα $4+1=5$, δηλαδή $16{,}5$. Εκατοστά: επόμενο το $6$ άρα $9+1=10$, που δίνει $16{,}50$. Χιλιοστά: επόμενο το $3$ άρα $16{,}496$.\par
\apantisi $16$, $16{,}5$, $16{,}50$, $16{,}496$.\par
\elegxos Το σφάλμα στα εκατοστά είναι $16{,}50-16{,}4963=0{,}0037<0{,}005$, δηλαδή μικρότερο από μισό εκατοστό. Σωστό. Το μηδενικό στο $16{,}50$ μένει για να δείχνει ότι μετρήσαμε μέχρι τα εκατοστά.

\paradeigma{Ένας δεκαδικός ανάμεσα}
\tixero Οι αριθμοί $3{,}14$ και $3{,}15$.\par
\tithelo Να βρω έναν δεκαδικό ανάμεσά τους.\par
\possindeo Γράφω τα δύο με τρία δεκαδικά: $3{,}140$ και $3{,}150$. Ανάμεσα είναι όλοι οι $3{,}141,\,3{,}142,\dots,3{,}149$.\par
\apantisi Για παράδειγμα ο $3{,}145$, που είναι ο μέσος τους.\par
\elegxos $3{,}14<3{,}145<3{,}15$. Και αφού μπορώ να ξαναχωρίσω το διάστημα, υπάρχουν άπειροι.

\tmima{Ασκήσεις}

\zoni{Το κατάλαβα}{δεκαδικά κλάσματα, ανάγνωση, σύγκριση}

\ex Γράψε ως δεκαδικούς: α) $\frac{7}{10}$ \quad β) $\frac{43}{100}$ \quad γ) $\frac{9}{1\,000}$ \quad δ) $\frac{2051}{1\,000}$ \quad ε) $\frac{38}{10\,000}$ \quad στ) $\frac{5}{100\,000}$.

\ex Γράψε ως δεκαδικά κλάσματα (με έναν αριθμητή και έναν παρονομαστή): α) $0{,}6$ \quad β) $0{,}07$ \quad γ) $3{,}25$ \quad δ) $0{,}0019$ \quad ε) $12{,}5$ \quad στ) $0{,}400$.

\ex Μετάτρεψε σε δεκαδικούς με διαίρεση ή με ισοδύναμο κλάσμα: α) $\frac{3}{4}$ \quad β) $\frac{7}{8}$ \quad γ) $\frac{9}{20}$ \quad δ) $\frac{11}{25}$ \quad ε) $\frac{5}{16}$ \quad στ) $\frac{3}{5}$.

\ex Γράψε με ψηφία: α) τρία ακέραια και είκοσι πέντε χιλιοστά \quad β) τετρακόσια εβδομήντα δύο εκατοστά \quad γ) δύο δέκατα \quad δ) δεκαπέντε δεκάκις χιλιοστά \quad ε) έξι ακέραια, πέντε εκατοστά και ένα εκατοντάκις χιλιοστό.

\ex Ποιο ψηφίο βρίσκεται στα εκατοστά και ποιο στα δεκάκις χιλιοστά του $24{,}0713$; Πώς διαβάζεται ο αριθμός;

\ex Ανάλυσε τους αριθμούς όπως στη θεωρία: α) $4{,}073$ \quad β) $12{,}408$ \quad γ) $0{,}0905$.

\ex Σωστό ή λάθος; α) $0{,}5=0{,}50$ \quad β) $0{,}5=0{,}05$ \quad γ) $3{,}40=3{,}4$ \quad δ) $0{,}0070=0{,}007$ \quad ε) Ο $12{,}5$ διαβάζεται «δώδεκα και πέντε» \quad στ) $\frac{1}{4}=0{,}4$ \quad ζ) $\frac{25}{100}=0{,}25$.

\ex Σύγκρινε τα ζευγάρια με τα σύμβολα $<,\,=,\,>$: α) $3{,}07$ και $3{,}7$ \quad β) $0{,}58$ και $0{,}6$ \quad γ) $12{,}09$ και $12{,}090$ \quad δ) $5{,}999$ και $6{,}001$ \quad ε) $0{,}3$ και $0{,}29$ \quad στ) $7{,}45$ και $7{,}5$.

\ex Στρογγυλοποίησε τους $7{,}486$ και $0{,}0951$: α) στη μονάδα \quad β) στα δέκατα \quad γ) στα εκατοστά.

\zoni{Το χρησιμοποιώ}{ο στίβος, η ζυγαριά, η πινακίδα}

\ex \textbf{Ο στίβος.} Στα 60 μ. οι χρόνοι ήταν: Νίκος $9{,}38$ s, Βασίλης $9{,}4$ s, Στέλιος $9{,}29$ s, Πολυξένη $9{,}35$ s, Αγγελική $9{,}3$ s, Ηλίας $9{,}41$ s. Γράψε τους μαθητές με τη σειρά που τερμάτισαν, από τον πρώτο.

\ex \textbf{Το άλμα.} Στο άλμα εις μήκος η Ζωή πήδηξε $3{,}08$ μ., η Μυρτώ $2{,}98$ μ., η Φανή $3{,}1$ μ. και η Πολυξένη $3{,}085$ μ. Γράψε τα μήκη από το μεγαλύτερο στο μικρότερο και πες ποια παιδιά πήραν το χρυσό και το ασημένιο μετάλλιο.

\ex \textbf{Η ζυγαριά.} Η Αγγελική ζυγίζει τα μήλα του παππού και βλέπει $2{,}35$ κιλά. Πόσα είναι στρογγυλοποιημένα στο κιλό και πόσα στα δέκατα του κιλού;

\ex \textbf{Η πινακίδα.} Στο βενζινάδικο η πινακίδα γράφει $1{,}879$ € το λίτρο. Ο θείος Λευτέρης λέει ότι «κοντά στα $1{,}9$», και η Πολυξένη ότι «κοντά στα $1{,}88$». Ποια στρογγυλοποίηση είναι ποια, και ποια είναι πιο κοντινή στην πραγματική τιμή;

\ex \textbf{Ανάμεσα.} Βρες τρεις διαφορετικούς δεκαδικούς ανάμεσα στους $2{,}4$ και $2{,}5$. Πόσοι ακόμη υπάρχουν;

\ex \textbf{Στην ευθεία.} Σχεδίασε ένα ευθύγραμμο τμήμα με μήκος $10$ εκ., που θα παριστάνει την απόσταση από το $0$ ως το $1$, και σημείωσε τους $0{,}3$, $0{,}75$, $0{,}9$ και $\frac{1}{2}$.

\ex \textbf{Τελειώνει ή όχι;} Κάνε τις διαιρέσεις (μέχρι να τελειώσουν ή να δεις ότι το ίδιο ψηφίο επαναλαμβάνεται): α) $\frac{3}{20}$ \quad β) $\frac{7}{30}$ \quad γ) $\frac{11}{50}$ \quad δ) $\frac{5}{12}$. Ποια τελειώνουν και ποια όχι;

\zoni{Το σκέφτομαι}{όταν το κόμμα ξεγελάει}

\ex Ο Βασίλης ισχυρίζεται ότι $0{,}9<0{,}12$, γιατί το $9$ είναι μικρότερο από το $12$. Εξήγησε το λάθος του και βρες ένα παράδειγμα με λεφτά ($0{,}90$ € και $0{,}12$ €) που το δείχνει.

\ex Πόσα δεκαδικά ψηφία έχει ο $\frac{7}{16}$; Πώς θα το προέβλεπες χωρίς να κάνεις τη διαίρεση, αν σκεφτείς ότι ο παρονομαστής $16$ είναι $2^4$ και ότι το $10^4=2^4\cdot5^4$;

\ex Με τα ψηφία $3$, $0$ και $8$ (το καθένα μία φορά στο δεκαδικό μέρος) και ακέραιο μέρος το $1$, γράψε τον μικρότερο και τον μεγαλύτερο δεκαδικό.

\ex α) Πόσοι δεκαδικοί, με το πολύ δύο δεκαδικά ψηφία, υπάρχουν ανάμεσα στο $2$ και στο $3$ (χωρίς αυτά); β) Πόσοι από αυτούς έχουν ακριβώς δύο δεκαδικά ψηφία (το τελευταίο όχι μηδέν);

\ex Ο Ηλίας στρογγυλοποιεί τον $2{,}449$ στα δέκατα με δύο βήματα: πρώτα στα εκατοστά ($2{,}45$) και μετά στα δέκατα ($2{,}5$). Ο Νίκος τον στρογγυλοποιεί κατευθείαν. Ποιος έχει δίκιο;

\ex Ο Στέβιν έγραφε τον $7{,}206$ ως $7_{(0)}\,2_{(1)}\,0_{(2)}\,6_{(3)}$. Πώς θα έγραφε τον $18{,}053$; Ποιος αριθμός είναι ο $8_{(0)}\,0_{(1)}\,4_{(2)}\,9_{(3)}$;

\zoni{Για το σπίτι}{μετρώ και καταγράφω}

\ex Μέτρησε το πλάτος του θρανίου σου με ένα μολύβι, με ένα χάρακα και με μετροταινία. Γράψε τα αποτελέσματα και εξήγησε ποια μέτρηση είναι η πιο ακριβής και γιατί χρειάζεται δεκαδικός αριθμός.

\ex Στο σούπερ μάρκετ βρες τρία προϊόντα που η τιμή ανά κιλό ή ανά λίτρο έχει τρία δεκαδικά ψηφία. Στρογγυλοποίησέ τις στο λεπτό του ευρώ και βάλε τες σε σειρά από την πιο ακριβή.
